\clearpage
\section{Refocusing for copying and allocation}
\label{sec:copying-intrinsic-refocusing}
A rooted state is reachable by search alone from the input of its own
underlying graph. Refocusing requires that a machine transition preserve this
property, even when the transition changes that graph. We now prove it for
every actual copying rewrite and every actual reference allocation, in arbitrary
surrounding graphs. The copying result includes an exact search-prefix length
law. The proof uses the complete successor certificates of
Section~\ref{sec:copying-successor-coherence}, not an assumed noninterference
property of copying.

\termdef{copying-intrinsic-refocusing}
\emph{Intrinsic copying refocusing} is preservation of search rootedness by the
search and copying branches of the existing combined relation. This partition
is exhaustive together with the operation-headed compute branch. Allocation
is then discharged within the latter branch. Full refocusing is proved
equivalent to refocusing for precisely the remaining operation-headed rewrites
other than reference creation, both at finite allowance and for the unbounded
machine. This is a reduction of outstanding obligations, not a proof that all
operations refocus. Section~\ref{sec:copying-full-refocusing} subsequently
discharges every remaining concrete operation case.

The paper correspondence is the copying branch of PDF Lemma~10.4
(HTML Lemma~9, \href{https://arxiv.org/html/1907.01257v5\#S10.SS3}{Section~10.3})
and the allocation branch of PDF Lemma~12.4 (HTML Lemma~13,
\href{https://arxiv.org/html/1907.01257v5\#S12.SS1}{Section~12.1}).
The prerequisites are operation-path reachability, PDF Lemma~C.4
(HTML Lemma~23), and the earlier-search-visit property, PDF Lemma~C.5
(HTML Lemma~24), in \href{https://arxiv.org/html/1907.01257v5\#A3}{Appendix~C}.
Our proof specializes the prefix-replay argument behind PDF Lemma~C.8
(HTML Lemma~27): copying uses its complete retained-occurrence description;
allocation uses replacement of a single operation edge. We do not yet claim
the full stationary-transition theorem for arbitrary copying graph fragments.

The dependency map for this milestone is:
\begin{enumerate}
\item Search rootedness implies operation-only reachability. A term-sharing
edge consequently blocks search from the retained copyable payload.
\item The original copy certificate provides retained vertices, edges and
ordered incidences in the complete reduct. The reachable prefix avoids the
only changed incidences, so it replays exactly.
\item Search history records earlier visits to operation sources, including
when control returns from an evaluated argument. It excludes revisiting a
source before an eventual rewrite normal form.
\item This history replays the prefix preceding a replaced operation head.
The already established allocation representation theorem identifies that
replacement with every actual allocation successor.
\item Recursive representation and padding preserve rootedness. The remaining
compute families are an explicit premise only in the full-run corollaries.
\end{enumerate}
The terms of Proposition~12.2, Figure~16's translation, the transition relation,
and all original stability and copyability conditions are unchanged.
The focused-graph correspondence, specimen locality, templates and final
equivalence chain remain open. Section~\ref{sec:copying-full-refocusing}
completes operation refocusing. The estimate
remains 55--70\% of mathematical effort, with low confidence; no percentage is
inferred from this section's 102 declarations.

\Needspace{12\baselineskip}
\subsection{Rootedness and operation paths}

\termdef{copying-search-rootedness}
For a graph $G$, an external input $i$, and a coordinate state $p$, write
$\operatorname{RootedFrom}(G,i,p)$ when
\[
 \exists n\in\mathbb N,\qquad (?,\operatorname{in}_i)
       \xrightarrow{\mathrm{search},n}_{G}p.
\]
The graph is fixed throughout this witness. In particular, a past execution
from a different graph is not a rootedness witness for the current state.

\lean{Copying/RootedSearch}{8}{9}{Search.RootedFrom}

\termdef{copying-operation-reachability}
\emph{Operation reachability} is generated by external inputs and by moving
from the unique source of an operation edge to one of its targets. The
operation-source lemma derives the singleton source list from the signature
and a membership hypothesis. This reachability predicate is an invariant
implied by rootedness, not an alternative definition of rootedness.

\lean{Copying/RootedSearch}{12}{16}{Search.OperationReachable}

\lean{Copying/RootedSearch}{19}{25}{Search.operation_source}

If a reachable vertex is a target of an edge with singleton source, that
source is reachable. An input cannot be an edge target; otherwise the last
operation in the reachability witness is the given incoming edge by target
uniqueness. The same argument shows that a target of a non-operation edge
cannot be operation reachable. Neither result assumes the surrounding graph
is acyclic.

\lean{Copying/RootedSearch}{28}{36}{Search.OperationReachable.parent}

\lean{Copying/RootedSearch}{39}{45}{Search.OperationReachable.not_nonoperation_target}

Every search interaction preserves operation reachability. Entering an
argument extends an operation path. Returning from an argument recovers its
operation parent, then either stays there or descends to the next argument.
Interactions which only change mode keep the same vertex; primitives with no
search interaction give impossible cases. Induction extends this statement
to runs, and starting at an input proves it for every rooted token.

\lean{Copying/RootedSearch}{48}{76}{Search.step_reachable}

\lean{Copying/RootedSearch}{78}{82}{Search.run_reachable}

\lean{Copying/RootedSearch}{84}{87}{Search.rooted_reachable}

The initial search state is rooted by the zero-length run. Appending one
search transition preserves rootedness. A recursive graph isomorphism preserves
and reflects it by the existing exact search-run transport, which fixes each
external input.

\lean{Copying/RootedSearch}{89}{90}{Search.rooted_initial}

\lean{Copying/RootedSearch}{92}{95}{Search.rooted_step}

\lean{Copying/RootedSearch}{97}{100}{Search.rooted_iso}

\Needspace{18\baselineskip}
For the closed interface $[\star]\Rightarrow[]$, \texttt{Rooted} chooses the
sole input. \texttt{UnboundedRooted} applies this predicate at the allowance
stored in an unbounded state. A state isomorphism maps the graph and token
together, so rootedness is invariant; every initial closed program is rooted.
The existing unbounded representation relation inherits the same equivalence
by its finite representative.

\lean{Copying/RootedSearch}{105}{105}{Combined.Rooted}

\lean{Copying/RootedSearch}{107}{107}{Combined.UnboundedRooted}

\lean{Copying/RootedSearch}{109}{111}{Combined.rooted_iso}

\lean{Copying/RootedSearch}{113}{114}{Combined.initial_rooted}

\lean{Copying/RootedSearch}{116}{117}{Combined.represents_rooted}

\Needspace{12\baselineskip}
\subsection{The retained payload is inaccessible to search}

A nonempty copy-resolution trace has a final sharing hop: induction on its
tail either returns the first hop or the final hop of the shorter trace. A
resolution from an active term contraction cannot have length zero, because
its terminal primitive would then also be that contraction, contradicting
uniqueness of the outgoing head. Thus it has a predecessor sharing edge.

\lean{Copying/CopySearchBarrier}{7}{11}{CopyTrace.last_hop}

\lean{Copying/CopySearchBarrier}{13}{21}{CopyTrace.resolves_predecessor}

\termdef{copying-search-barrier}
A \emph{copying search barrier} is the obstruction supplied by a sharing edge
immediately before the payload. Forwarding labels are contractions or
weakenings, never operations. The final hop targets the payload's singleton
source, so the non-operation-target lemma excludes operation reachability
there. This argument uses the actual incidence relation in the whole host.

\lean{Copying/CopySearchBarrier}{23}{27}{CopyTrace.forward_not_operation}

\lean{Copying/CopySearchBarrier}{29}{42}{CopyTrace.resolved_source_unreachable}

A successor match includes an occurrence of the complete payload and its
resolution trace. Mapping its singleton head source into the host therefore
proves that the payload input is unreachable by search. For a primitive
payload its tail is its head; for a lambda payload its tail is a box, which
cannot satisfy the primitive-label premise. Consequently no reachable
primitive search interaction can use the retained tail's source. This is the
specific fact needed when copying changes the tail's targets.

\lean{Copying/CopySearchBarrier}{44}{48}{CopyRule.SuccessorMatch.payload_input_unreachable}

\lean{Copying/CopySearchBarrier}{51}{58}{CopyRule.SuccessorMatch.primitive_tail_eq_head}

\lean{Copying/CopySearchBarrier}{61}{68}{CopyRule.SuccessorMatch.tail_source_unreachable}

\Needspace{12\baselineskip}
\subsection{Retained occurrences inside the actual reduct}

\termdef{copying-retained-occurrence}
A \emph{retained occurrence} identifies an old source vertex or edge inside
the complete copying successor. The certificate's realization covers every
old occurrence. Choose a preimage for each old internal vertex and edge, and
extend the vertex map to ports while fixing the external interface. These
choices use ordinary classical choice; they do not assert that the whole old
graph is an unchanged subgraph.

\lean{Copying/CopyRetainedOccurrence}{10}{11}{CopyRule.SuccessorCertificate.retainedVertex}

\lean{Copying/CopyRetainedOccurrence}{12}{13}{CopyRule.SuccessorCertificate.retainedEdge}

\lean{Copying/CopyRetainedOccurrence}{14}{14}{CopyRule.SuccessorCertificate.retainedPort}

The choice specifications say that the chosen vertex and edge realize the
corresponding old occurrence. The port equation follows by its three cases:
internal, input and output. Applying the realization to two equal retained
ports gives equality of their old occurrences, hence injectivity of the port
map.

\lean{Copying/CopyRetainedOccurrence}{16}{17}{CopyRule.SuccessorCertificate.retained_vertex}

\lean{Copying/CopyRetainedOccurrence}{18}{19}{CopyRule.SuccessorCertificate.retained_edge}

\lean{Copying/CopyRetainedOccurrence}{20}{21}{CopyRule.SuccessorCertificate.retained_port}

\lean{Copying/CopyRetainedOccurrence}{23}{25}{CopyRule.SuccessorCertificate.retained_port_injective}

Every retained primitive has the same label, by the certificate's recursive
label comparison. Every source list other than that of the active contraction
is the old ordered source list mapped through retained ports. Every target
list other than that of the payload tail obeys the analogous law. Apply the
injective realization map to both sides and use the source-relative
presentation equations. The two exceptions are essential: copying really
changes those attachments. Finally, the certificate's token equation
identifies the retained request vertex with the resumed vertex.

\lean{Copying/CopyRetainedOccurrence}{27}{31}{CopyRule.SuccessorCertificate.retained_label}

\lean{Copying/CopyRetainedOccurrence}{33}{39}{CopyRule.SuccessorCertificate.retained_source}

\lean{Copying/CopyRetainedOccurrence}{41}{47}{CopyRule.SuccessorCertificate.retained_target}

\lean{Copying/CopyRetainedOccurrence}{49}{50}{CopyRule.SuccessorCertificate.retained_token}

\Needspace{12\baselineskip}
\subsection{Exact replay of the copying prefix}

Pure search has no outgoing transition from rewrite mode. Therefore a search
run ending outside rewrite mode cannot start in it: normality would force
its endpoint to equal its source. A separate induction splits any run of
length $n+1$ into a prefix of length $n$ and its last interaction.

\lean{Copying/CopySearchReplay}{8}{17}{Search.run_source_not_rewrite}

\lean{Copying/CopySearchReplay}{20}{29}{Search.last_step}

An entering interaction which produces rewrite mode puts the token at the
vertex it was given; so does an operation-resume interaction. This follows
by inspecting the actual \texttt{enter} and \texttt{resume} equations. If the
resulting vertex has a term-contraction head, the resume case is impossible:
it would give the same vertex an operation head. Hence the immediate
predecessor of a copying request is search at that same vertex.

\lean{Copying/CopySearchReplay}{31}{44}{Search.enter_rewrite_vertex}

\lean{Copying/CopySearchReplay}{46}{51}{Search.resume_rewrite_vertex}

\lean{Copying/CopySearchReplay}{54}{65}{Search.contraction_predecessor}

The zero-length run cannot end in rewrite mode. Splitting a rooted copying
request at its last interaction and applying the predecessor result gives
$n=m+1$ and an exact $m$-step run ending at search on the request vertex.
The last step is precisely the copying trigger.

\lean{Copying/CopySearchReplay}{68}{83}{Search.copying_prefix}

Take a reachable search interaction whose successor is not in rewrite
mode. It cannot use the active contraction: entering that head would request
a rewrite. Nor can it use the payload tail: the search barrier excludes its
source, and in a backtracking case the reachable argument has a reachable
parent. All required ordered incidences are therefore preserved by the
retained maps. The existing naturality equations for entering and resuming
then construct the same interaction in the actual successor. Induction
replays a whole prefix; its intermediate states cannot be in rewrite mode
because its endpoint is not.

\lean{Copying/CopySearchReplay}{86}{114}{CopyRule.SuccessorCertificate.replay_step}

\lean{Copying/CopySearchReplay}{117}{126}{CopyRule.SuccessorCertificate.replay_prefix}

\termdef{copying-prefix-refocusing-cost}
The \emph{copying refocusing cost} is an exact comparison between searches on
two different graphs. If a search on $G$ reaches the request in $n$ steps,
then, for an actual copy successor $K$, there is an $m$ such that
\[
 n=m+1,\qquad (?,\operatorname{in}_i)
       \xrightarrow{\mathrm{search},m}_{K}(?,w).
\]
Replay the pre-trigger prefix and use the retained-token equation. Forgetting
$m$ gives rootedness of the successor. This statement counts search steps;
it does not identify the graph-changing copy transition with a search step.

\lean{Copying/CopySearchReplay}{128}{135}{CopyRule.SuccessorCertificate.rooted_cost}

\lean{Copying/CopySearchReplay}{137}{140}{CopyRule.SuccessorCertificate.rooted}

\Needspace{12\baselineskip}
\subsection{Discharging the intrinsic machine cases}

The public copying theorems take an actual contextual rewrite at a
term-contraction head. The existing extraction theorem supplies its complete
successor certificate; the preceding results then give rootedness and the
exact prefix-length law. No certificate or noninterference premise is left
for the caller to prove.

\lean{Copying/IntrinsicRefocusing}{7}{11}{Combined.copying_rooted}

\lean{Copying/IntrinsicRefocusing}{14}{19}{Combined.copying_rooted_cost}

The intrinsic relation consists of unchanged search transitions and actual
rewrites at term-contraction heads. The compute relation consists of actual
rewrites at operation heads. Both embed directly into the original combined
relation. They cannot both start at the same state: search has no rewrite
source, and a contraction head cannot also be an operation head. Canonical
source recognition shows that every combined step belongs to one of these
two relations, so this is an exhaustive partition of the existing machine.

\lean{Copying/IntrinsicRefocusing}{22}{25}{Combined.IntrinsicStep}

\lean{Copying/IntrinsicRefocusing}{28}{30}{Combined.ComputeStep}

\lean{Copying/IntrinsicRefocusing}{32}{35}{Combined.IntrinsicStep.combined}

\lean{Copying/IntrinsicRefocusing}{37}{38}{Combined.ComputeStep.combined}

\lean{Copying/IntrinsicRefocusing}{40}{45}{Combined.intrinsic_compute_disjoint}

\lean{Copying/IntrinsicRefocusing}{48}{60}{Combined.step_iff_intrinsic_or_compute}

Refocusing for a relation means that a rooted source has a rooted successor.
The intrinsic relation refocuses: the search case appends a search step, and
the copy case uses the newly proved theorem. Hence the whole finite machine
is refocusing if and only if its compute relation is refocusing. The forward
implication restricts a full-machine proof; the reverse implication uses
the exhaustive partition.

\lean{Copying/IntrinsicRefocusing}{62}{63}{Combined.Refocusing}

\lean{Copying/IntrinsicRefocusing}{65}{69}{Combined.intrinsic_refocusing}

\lean{Copying/IntrinsicRefocusing}{72}{78}{Combined.refocusing_iff_compute}

Lift each of the two relations to the unbounded state by keeping the same
finite allowance for a step. Their disjunction is again exactly the actual
unbounded relation. Unbounded refocusing is preservation of the unbounded
rooted predicate. The intrinsic theorem lifts without additional premises,
and the full unbounded obligation reduces equivalently to the compute
obligation. No global upper bound on future box depth is asserted.

\lean{Copying/IntrinsicRefocusing}{80}{81}{Combined.UnboundedIntrinsicStep}

\lean{Copying/IntrinsicRefocusing}{82}{83}{Combined.UnboundedComputeStep}

\lean{Copying/IntrinsicRefocusing}{85}{95}{Combined.unbounded_step_iff_intrinsic_or_compute}

\lean{Copying/IntrinsicRefocusing}{97}{98}{Combined.UnboundedRefocusing}

\lean{Copying/IntrinsicRefocusing}{100}{102}{Combined.unbounded_intrinsic_refocusing}

\lean{Copying/IntrinsicRefocusing}{104}{112}{Combined.unbounded_refocusing_iff_compute}

Once compute refocusing is supplied, induction over an actual finite run
preserves rootedness at its endpoint. The compute-refocusing hypothesis is explicit in this corollary.
Section~\ref{sec:copying-full-refocusing} supplies it for the concrete machine
and proves the unconditional full-run theorem.

\lean{Copying/IntrinsicRefocusing}{115}{120}{Combined.run_rooted}

\Needspace{12\baselineskip}
\subsection{Exact search history, including backtracking}

\termdef{copying-search-visit}
A \emph{search visit} to $v$ in a run ending at $p$ with length $n$ is an exact
factorization $n=a+b$ into a run from the input to $(?,v)$ followed by a run
to $p$. \emph{Search history stacks} are operation paths from the chosen input
whose operation sources satisfy a predicate $P$ recording earlier visits.
They are logical witnesses reconstructed from a run, not an added runtime
stack.

\lean{Copying/SearchHistory}{8}{10}{Search.Visit}

\lean{Copying/SearchHistory}{12}{16}{Search.HistoryStack}

Replacing $P$ by a pointwise consequence $Q$ preserves a history stack by
induction. From a stack at an edge target, recover its parent stack and that
parent's predicate: inputs are not targets, and uniqueness of the incoming
edge identifies the final operation. A recorded visit extends across one
additional step by extending its suffix; the current search vertex is visited
with suffix length zero.

\lean{Copying/SearchHistory}{18}{22}{Search.HistoryStack.mono}

\lean{Copying/SearchHistory}{24}{31}{Search.HistoryStack.parent}

\lean{Copying/SearchHistory}{33}{35}{Search.Visit.step}

\lean{Copying/SearchHistory}{37}{38}{Search.Visit.current}

A full history consists of a stack annotated by exact visits in the current
run, together with a visit to its current token vertex. For one search step,
extend every recorded suffix. Entering a fresh argument adds its operation
source to the stack; resuming recovers that source from the parent theorem,
then either returns to it or records the next argument. Mode-only changes
retain the vertex's earlier visit. This proves history preservation for every
search rule, including backtracking.

\lean{Copying/SearchHistory}{40}{41}{Search.History}

\lean{Copying/SearchHistory}{43}{76}{Search.history_step}

Induction extends history across a run, adding its length to the existing
history length. The initial input has the empty stack and its zero-length
visit, so every rooted run carries a full history. In particular, if its final
mode is not search, the visit to its final vertex has a strictly positive
suffix: a zero suffix would identify the endpoint with search mode. This
proves the exact-run version of PDF Lemma~C.5.

\lean{Copying/SearchHistory}{78}{83}{Search.history_run}

\lean{Copying/SearchHistory}{85}{87}{Search.history}

\lean{Copying/SearchHistory}{90}{98}{Search.revisit}

\Needspace{12\baselineskip}
\subsection{Single operation replacement and allocation}

A prior search visit cannot recur with a positive intervening run before an
eventual rewrite endpoint. Otherwise concatenate the intervening cycle with
the suffix to obtain two search runs from the same state to a rewrite normal
form with unequal lengths. Determinism and normality imply equal lengths,
contradicting positivity. Only the already proved determinism of search is
used here.

\lean{Copying/HeadReplacementRefocusing}{9}{21}{Search.Visit.no_earlier_cycle}

Replace one edge $e$ by an arbitrary graph $A$ of the same typed interface.
A search interaction which does not use $e$ transports along the retained
host-port map. Its other edge is retained with the same primitive label and
mapped ordered source and target lists; the search equations commute with
this map. The premise specifies both ways an interaction can use an edge:
entering one of its sources or returning at an ordered target.

\lean{Copying/HeadReplacementRefocusing}{24}{45}{Search.replace_step}

Consider a run prefix ending at search on the source $v$ of $e$, followed
by a run to a rewrite state. No earlier interaction in that prefix uses $e$.
An entering interaction would already be a visit to $v$; a return at a target
of $e$ would imply such an earlier visit by search history. Either contradicts
the no-earlier-cycle result. Induction therefore transports the entire prefix
through replacement, with exactly the same length.

\lean{Copying/HeadReplacementRefocusing}{48}{71}{Search.replace_prefix}

\termdef{copying-head-replacement-refocusing}
\emph{Head replacement refocusing} says that if a rewrite request is rooted
at the source of an operation, replacing just that edge by any graph with its
interface leaves search at the retained source rooted. The earlier-visit
theorem supplies the pre-head prefix; the preceding replay theorem transports
it. This special theorem needs no one-way premise on the replacement graph.
It does not authorize deletion or rewiring of other host edges, and therefore
does not by itself cover all compute rules.

\lean{Copying/HeadReplacementRefocusing}{74}{82}{Search.replacement_rooted}

Recursive state isomorphism preserves rootedness from any specified input,
not just a closed program's input. The graph transport carries the witness,
and the state-isomorphism equation identifies its endpoint token. For an
actual allocation rewrite, the earlier canonical-successor theorem supplies
an isomorphism from replacement of the original reference head by its fresh
cell prefix to the actual successor. Apply head replacement refocusing and
then this isomorphism. The allocation rule still requires the same stable
argument as before; the proof never removes that premise.

\lean{Copying/AllocationRefocusing}{7}{10}{Combined.StateIso.rootedFrom}

\lean{Copying/AllocationRefocusing}{13}{17}{Combined.allocation_rooted}

The remaining compute relation contains actual rewrites at operation heads
whose operation is not reference creation. It embeds in the compute relation.
To prove full refocusing from refocusing of this remaining relation, partition
the operation cases by equality with reference creation: use the allocation
theorem in that case and the supplied premise otherwise. Conversely, any
full-machine proof restricts to the remaining relation. Thus this reduction
is an equivalence, not merely a sufficient-condition theorem.

\lean{Copying/AllocationRefocusing}{20}{23}{Combined.RemainingComputeStep}

\lean{Copying/AllocationRefocusing}{25}{27}{Combined.RemainingComputeStep.compute}

\lean{Copying/AllocationRefocusing}{29}{39}{Combined.refocusing_iff_remaining_compute}

Lift the remaining relation to unbounded states and repeat the same
argument at each step's finite allowance. Finally, induction preserves
rootedness along any finite combined run once refocusing of this remaining
relation has been proved. This version retains an explicit premise covering dereference, assignment,
equality, application and arithmetic. Section~\ref{sec:copying-full-refocusing}
now discharges it for the concrete machine, without assuming termination
of any continuation.

\lean{Copying/AllocationRefocusing}{41}{43}{Combined.UnboundedRemainingComputeStep}

\lean{Copying/AllocationRefocusing}{45}{57}{Combined.unbounded_refocusing_iff_remaining_compute}

\lean{Copying/AllocationRefocusing}{60}{65}{Combined.run_rooted_of_remaining}

\Needspace{12\baselineskip}
\subsection{Representation laws and semantic regressions}

A search transport between graphs with one term input must map that input
to itself: it preserves the input predicate, and there is only one input
index. Exact transport and reflection of runs therefore preserve and reflect
rootedness. Applying this result to the already proved concrete padding
transport establishes rootedness invariance at every larger allowance. This
uses concrete run transport, not an inference from determinism alone.

\lean{Copying/RootedPadding}{8}{15}{Search.Transport.sole_input}

\lean{Copying/RootedPadding}{17}{22}{Search.Transport.rooted_iff}

\lean{Copying/RootedPadding}{24}{26}{Search.rooted_padTo_iff}

The retained payload input is not rooted in any focus mode or from any
external input, since every rooted token would be operation reachable. For
closed states, any combined successor of a rooted contraction request is
rooted: a search successor from rewrite mode is impossible, leaving the
actual copying theorem.

\lean{Copying/RefocusingExamples}{8}{16}{CopyRule.SuccessorMatch.payload_input_not_rooted}

\lean{Copying/RefocusingExamples}{18}{23}{Combined.copying_step_rooted}

The test graph uses a canonical one-input contraction tree and a primitive
reference-allocation operation as an enclosing host. The two payload graphs
are the original \texttt{CopyRule.before} and \texttt{CopyRule.after} applied
to a constant. The host first enters its argument, where the copying request
is now an internal token position. Entering the host followed by the copy
trigger gives an exact two-step rootedness witness.

\lean{Copying/RefocusingExamples}{27}{27}{Combined.RefocusingExamples.copyTree}

\lean{Copying/RefocusingExamples}{28}{28}{Combined.RefocusingExamples.copyHost}

\lean{Copying/RefocusingExamples}{29}{30}{Combined.RefocusingExamples.copiedConstant}

\lean{Copying/RefocusingExamples}{31}{32}{Combined.RefocusingExamples.copiedConstantAfter}

\lean{Copying/RefocusingExamples}{34}{39}{Combined.RefocusingExamples.host_enters_copy}

\lean{Copying/RefocusingExamples}{42}{45}{Combined.RefocusingExamples.host_copy_request_rooted}

The actual copying step under this host preserves rootedness by the general
theorem. The returned payload at the host's internal argument is not a final
state; the outer reference operation is still pending. This regression checks
that refocusing and finality are kept distinct at an internal token, rather
than testing only a copy request at the external root.

\lean{Copying/RefocusingExamples}{48}{53}{Combined.RefocusingExamples.host_copy_result_rooted}

\lean{Copying/RefocusingExamples}{55}{57}{Combined.RefocusingExamples.host_copy_return_not_final}

For the unchanged translation of $\mathtt{ref}\ n$, the existing entering,
value and resume interactions give a three-step search witness to the
allocation request. Every actual combined successor is rooted by allocation
refocusing; no particular representative of the allocated graph is assumed.
The unchanged translated addition example has its established five-step
search witness, and its rootedness holds under every finite padding by the
proved padding equivalence.

\lean{Copying/RefocusingExamples}{60}{63}{Combined.RefocusingExamples.translated_reference_request_rooted}

\lean{Copying/RefocusingExamples}{66}{72}{Combined.RefocusingExamples.translated_reference_result_rooted}

\lean{Copying/RefocusingExamples}{75}{76}{Combined.RefocusingExamples.translated_addition_request_rooted}

\lean{Copying/RefocusingExamples}{78}{81}{Combined.RefocusingExamples.translated_addition_rooted_padding}

All 102 declarations above are included in the axiom audit. The full build
passes 446 jobs, and 2305 axiom checks pass: 6 abstract, 477 graph and 1822
copying declarations. Only \texttt{propext}, \texttt{Classical.choice} and
\texttt{Quot.sound} occur. Appendix~B.1 removes the discharged copying and
allocation refocusing sub-obligations. Section~\ref{sec:copying-full-refocusing}
now removes the other operation-refocusing cases as well. Specimen, template
and final-assembly obligations remain explicit.
